\chapter{Design}
\label{chap:design}

% USER INFORMATION TO PROFILE
% we established research gap
    % but still want to justify design decisions
% this is where the bias discussion comes in:
    % can cite my interesting paper where people prefer biased recs in the context of college courses
    % gender is established as a bias, how does that factor in?
    % parental education bg, SES, parental influence are also factors
    % i could make the argument that I wanted my RS to be "colour blind" to this
% don't want to overcomplicate and over-engineer the amount of preferences to ask a user



% owen: ok to mention AI as a starting point,
% given that I had a lot of manual labour on top of it with my learned knowledge regarding the subject matter


\section{Relevant information to question}
\label{sec:relevant}

% we must ask ourselves the question (which is a research objective): 
% what information is relevant to question a user in order to make effective, relevant recs for a CCRS?




% talk about thinking about whether to use LLM's or not

% owen: talk about manual labour




Interviews with expert guidance counsellors were conducted as part of this research, and a qualitative analysis of these interviews are provided in Section \ref{sec:interviews}.

\subsection{Interviews with guidance counsellors}
\label{sec:interviews}

\section{Problem Formulation}
\label{sec:ProblemFormulation}

This section should provide the reader with an overall description of the problem that will be addressed in the dissertation. In contrast to a generic discussion of the dissertation topic in the Introduction chapter, this section should provide a detailed discussion of the problem that has been identified based on the existing work that has been discussed in the preceeding chapter. 

In some dissertations, it may make sense to convert this section into a short chapter of its own which follows the discussion of the existing work and preceeds the discussion of the work of the dissertation.

\subsection{Identified Challenges}
This section should present a short description of the gaps in the existing work and the relationship of these gaps to the work described in this dissertation.

\subsection{Proposed Work}
This section should provide a thorough description of the problem and an overview of the work proposed to address the problem.


\section{Overview of the Design}
\label{sec:OverviewOfDesign}
A description of the approach that addresses the problem identified above.


\section{Summary}
\label{sec:SummaryDesign}

Every chapter aside from the first and last chapter should conclude with a summary. 